\scanpage{114}
\deleted{B) Cas général. On se réduit au}

Si $F$ est quelconque, on écrit que $F$ est un sous-\deleted{quotient d'un \inserted{\unsure{ct tordu|ct. tordu}} quotient \inserted{\unsure{ct tordu}} d'un $R^df_*(Q_{lY})$}, où $f:Y\longrightarrow X$ morphisme projectif. [Utilisant Lefschetz \deleted{et 2.1 de \unsure{$f$, $f'$} \unsure{is. à une section de $f$}, faible, on peut supposer que les fibres de $f$ \unsure{sont de dim $j$}]} \deleted{\unsure{fort|fortement}} si $j<d=\dim$ relative de $Y/X$, et remplaçant $Y$ par $Y\times_X\mathbb P_X^{j-d}$ si $j>d$, on peut supposer $j=d$]
\margin{\deleted{\unsure{par Lefschetz fort.|Lefschetz fort.}}}
% slanted
\margin{voir si cette réduction est vraiment utile\ldots}
% slanted

La suite spectrale
\[
 R^\alpha_!(fg)(Q_l)\Longleftarrow E_2^{pq}=R^p_!f(R^q_!g(Q_{lX}))
\]
dégénère \uline{par Deligne}, donc
\[
\begin{aligned}
 R^i_!f(F)=E^{ij}_2\quad &\text{est }\simeq\text{ quotient \inserted{c$^{\mathrm t}$. tordu} d'un ss-faisceau ct}\\
 &\text{tordu d'un }R^{i+\unsure{j|g}}_!g(Q_{lY}).
\end{aligned}
\]
\margin{\uuline{utiliser Lefschetz fort}}
% slanted
\margin{[N.B. Voir si Deligne ne peut se remplacer ici par \uline{factorisation d'}\uline{Artin} sur $X\longrightarrow S$, pour revenir au cas de $n=1$\ldots]}
% slanted
La conclusion en résulte, en utilisant le cas précédent pour $Y\longrightarrow S$, de dim relative $\leq n+j$ \ldots

\uline{Corollaire} 5.2. Si la partie \uuline{\textquotedblleft semi-simplicité\textquotedblright} des conjectures de Tate est vérifiée, la conclusion de 5.1. est vraie dès que l'on suppose seulement que $F$ est eff$^{\mathrm t}$. admissible de poids $j$. Par suite, si on suppose $F$ eff$^{\mathrm t}$. admissible de poids $\leq\deleted{\unsure{q}}\inserted{j}$, on a un énoncé correspondant, \deleted{\unsure{De ce}} $R^if_!(F)$ est \inserted{eff. admissible,} \deleted{de} poids $\leq\deleted{\unsure{\operatorname{Inf}(i+j,2n-i+j)}}\unsure{i+j}$, \inserted{niveau $\leq\operatorname{Inf}(i+j,2n-i+j)$.} Si $i\geq n$, on trouve des poids entre $2(i-n)$ et $i+j$, correspondant aux niveaux $0,1,\ldots,\unsure{2-i+j|2n-i+j}$.
\margin{semi-simplicité utilisée!}
% slanted; the note is circled.

(On se ramène en effet \inserted{par dévissages et déf.} au cas où $F$ est \uline{spécial} de poids $j$, et par th. de chgt. de base propre au cas où $S$ de type fini sur Spec $\mathbb Z$, et \unsure{ainsi de suite}\deleted{\unsure{au cas où}})

\uline{Cor.} 5.3 (sous conj. type Tate) \deleted{\unsure{on a aussi}} Si $F$ est eff. admissible, Alors $R^i_!f(F)$ est eff$^{\mathrm t}$. admissible.
\margin{Alors par semi-simplicité $F$ est $=$ str$^{\mathrm t}$ spécial}
% slanted
% archive sheet number: 113
